\documentclass[]{article}
\usepackage[utf8]{inputenc}
\usepackage{multicol, amsmath, amssymb, hyperref, amsthm}


\title{Year 9 Worksheets}
\author{C C Barnes }
\date{September 2022}

\pagenumbering{gobble}
\usepackage{geometry}

\begin{document}

\maketitle
\begin{abstract}
    
\end{abstract}

\newpage
\section{Straight Line Graphs}

\newpage

\noindent \textbf{Perpendicular Lines}

\medskip

\hrule

\medskip

\begin{enumerate}
    \item Find the equation of the line passing through \((1,4)\) perpendicular to the line \(y=3x-1\). Give your answer in the form \(ax+by+c=0\) where \(a,b,c \in \mathbb{Z}\).
    \item Find the equation of the line perpendicular to \(2x+3y+5=0\) passing through \((-1,3)\). Give your answer in the form \(ax+by+c=0\) where \(a,b,c \in \mathbb{Z}\).
\end{enumerate}

\noindent \textbf{Extension}

\begin{enumerate}
    \item Line \(A\) passes through the points \((-3, -1)\) and \((-1, 9)\). Line \(B\) passes through the points \((-2, 1)\) and \((k, 4)\). Line \(A\) and Line \(B\) are perpendicular. Find the value of \(k\). 
\end{enumerate}

\newpage

\noindent \textbf{Answers (Extension only!)}

\medskip

\hrule

\medskip

\noindent Let \(m_A, m_B\) be the gradients of lines \(A\) and \(B\) respectively. Then
\[
m_A = \frac{9--1}{-1--3} = \frac{10}{2} = 5.
\]
Since \(B\) is perpendicular to \(A\),
\[
m_B = - \frac{1}{m_A} = -\frac{1}{5}.
\]
Thus \(B\) has equation
\[
y=-\frac{1}{5}x + c.
\]
Since \(B\) passes through \((-2,1)\),
\begin{align*}
    1 &= -\frac{1}{5}(-2) + c \\
    \implies  c &= \frac{3}{5} \\
    \implies y &= -\frac{1}{5}x + \frac{3}{5}.
\end{align*}
Thus when \(y=4\),
\begin{align*}
    4 &= -\frac{1}{5}x + \frac{3}{5} \\
    \implies x &= -17 \\
    \implies k &= -17.
\end{align*}

\noindent Can you find another method for solving this?

\newpage

\noindent \textbf{Gradient and Intercept}

\medskip
\hrule
\medskip

\begin{enumerate}
    \item In each of the following identify the gradient and intercept. Your answers need not be numerical and may include variables.
    \begin{enumerate}
        \item \(y=3x-4\)
        \item \(y+2=-3x\)
        \item \(y-7=3(x-2)\)
        \item \(3x+2y+1=0\)
        \item \(y-y_1=m(x-x_1)\)
        \item \(x+\pi=3y-2\)
        \item \(2x+3y + b= ax - 4\)
        \item \(\frac{x}{a}+\frac{y}{b}=1\)
    \end{enumerate}

    \item In each case, calculate the gradient between the two points.

    \begin{enumerate}
        \item  \((3,2) \text{ and } (1,7)\)
        \item \((-4, -2)\) and \((6,-1)\)
        \item \((\frac{1}{2}, 5)\) and \((3, \frac{2}{3})\)
        \item \((-\frac{2}{5}, \frac{1}{4})\) and \((\frac{1}{2}, -\frac{7}{8})\)
        \item \((x_1,y_1)\) and \((x_2,y_2)\)
        \item \((3a, 7a)\) and \((2a, -3a)\)
        \item \((2a-3,-2)\) and \((5a-3,7)\)
        \item \( (a,2a^2) \) and \((5,10a)\)
    \end{enumerate}

\end{enumerate}

\newpage

\section{Parabolas}

\newpage

\noindent \textbf{Function Notation}

\medskip

\hrule

\medskip

\noindent \textbf{Extension}

\begin{enumerate}
    \item Suppose we have two functions, \(f(x)=x^2+6x+1\) and \(g(x)=\frac{1+x}{2-x}\). What is \(f(g(3))\)?
    \item What is \(g(f(x))\)?
    \medskip
    \newline
    \textit{Aside: This process of combining functions in this way is known as composition of functions and we can denote the resulting function as \(f \circ g\). }
    \item  Pick any two functions \(f\) and \(g\). Is it true that \(f\circ g=g\circ f\)? This property is known as \textit{commutativity}.
    \item Suppose \(f(x)=\pi\) and \(g(x)=x+4\). Evaluate both \((f \circ g)(x)\) and \((g \circ f)(x)\).
\end{enumerate}

\newpage

\noindent \textbf{Answers}

\medskip

\hrule

\medskip

\begin{enumerate}
    \item \begin{enumerate}
        \item \(f(g(3))=-7\)
        \item \(g(f(x)) = \frac{x^2+6x+2}{1-6x-x^2}\)
    \end{enumerate}
    \item No. For example, let \(f(x)=x^2, g(x)=x+1\). Then \\
    \((f\circ g)(x) = x^2+2x+1\) but \\
    \((g\circ f)(x) = x^2 +1\).
    \item \((f\circ g)(x) = \pi\) \\
    \((g\circ f)(x) = \pi +4.\)
\end{enumerate}

\newpage

\section{Further Factorisation}

\newpage

\noindent \textbf{Further Factorisation}

\medskip

\hrule

\medskip

\noindent Factorise the following expressions:

\begin{enumerate}

    \item $7(x-2) + 3(x-2)$
    \item $5(y-6) - 7(y-6)$
    \item $12(z-4) - 7(4-z)$
    \item $16(2x-4) + 5(x-2)$
    \item $3(x-2)^3 - 2(x-2)$
    \item $(x^2-4x) - 4(x-4)$
    \item $y^3(7-y) + 3y(7-y)$
    \item $5z^2(z-3)^2 - 5(3z-z^2)^2 $

\end{enumerate}

\newpage

\noindent \textbf{Answers}

\medskip

\hrule

\medskip

\begin{enumerate}
    \item $10(x-2)$
    \item $-2(y-6)$ or $2(6-y)$ for +1 style point.
    \item $19(z-4)$
    \item $37(x-2)$
    \item $(x^2+4x+2)(x-2)$ \\
    \\
    For the record, the quadratic $x^2+4x+2$ can be factorised but it doesn't have integer roots. It's $(x+2-\sqrt{2})(x+2+\sqrt{2})$ for anyone who's interested... Go ahead and check by expanding it out!
    \item $(x-4)^2$
    \item $y(7-y)(y^2+3)$
    \item 0
\end{enumerate}

\newpage






\end{document}
